\chapter{Efficient Updates of the KKT System}
\label{chap:updates}

\vspace*{\fill}

Every iteration of the active-set method changes the working set by exactly one index: a
blocking variable is added, or a variable with a negative multiplier is released. The
previous chapter showed that everything the method needs is determined by the extended
ordering and by a factorization of the effective reduced Hessian $\Heff$. This chapter shows
that a change of the working set changes $\Heff$ only by a \textbf{low-rank modification}, so
that its $LDL^\top$ factorization can be updated in place instead of being recomputed.

\medskip

We first introduce a \textbf{fixed-length representation}, in which the factorization has a
fixed dimension and inactive indices are represented by identity rows. We then derive the
updates for the three ways of \textbf{expanding} the working set and the two ways of
\textbf{reducing} it, and translate them into row additions, row deletions and rank-one
modifications of the factorization. Row deletions leave structural zeros in the factor; we
analyse this effect and the \textbf{periodic refactorization} that controls it.

\medskip

Finally, we show that the state of the method can be carried over from one column generation
iteration to the next. A new column and a new stability center leave the factorization
unchanged, which gives a \textbf{hot start} of every master problem. We also analyse the
center-only iterations of \cref{sec:center-update} in this light.

\vspace*{\fill}

\newpage


\section{Fixed-Length Representation}
\label{sec:fixed-length}

The number of nullspace columns $|\Nx|$ changes in almost every iteration, and the number of
columns $n$ grows during column generation. Resizing a sparse factorization at every change
would be expensive and complicated. Instead, the factorized matrix is given a fixed dimension.

\begin{definition}[Fixed-length representation] \label{def:fixed-length}
Let $\nmax\ge n$ be a fixed \textbf{column capacity}, the maximal number of columns that column
generation may generate. The \textbf{embedded effective reduced Hessian} is the
$\nmax\times\nmax$ matrix
\[
\bigl(\Heff^{\mathrm{emb}}\bigr)_{kj} =
\begin{cases}
(\Heff)_{kj} = u_k^\top M_\rho u_j, & k,j\in\Nx,\\
1, & k = j\notin\Nx,\\
0, & \text{otherwise,}
\end{cases}
\]
with $u_j = a_j - a_{\rep(j)}$ and the row mask $M_\rho$ of~\eqref{eq:implicit}. Index $j$ of
the embedded matrix is column $j$ of the master problem. Indices of columns that are not in
$\Nx$, including columns that have not been generated yet, carry an identity row and column.
\end{definition}

In this representation, a column entering $\Nx$ replaces an identity row by a row of $\Heff$,
and a column leaving $\Nx$ is replaced by an identity row (\cref{fig:fixed-length}). The
dimension never changes. The surpluses do not appear at all: by \cref{thm:decoupling} they
have no index in the factorized matrix, so the inequality rows do not increase the size of
the factorization.

\begin{figure}[ht]
\centering
\begin{tikzpicture}[x=0.32cm, y=0.32cm]
    % active indices (row b of the matrix is drawn at height 16-b .. 17-b)
    \foreach \a in {2,5,6,10,13} {
        \foreach \b in {2,5,6,10,13} {
            \fill[cblue!60] ({\a-1},{16-\b}) rectangle ({\a},{17-\b});
        }
    }
    % identity on inactive
    \foreach \k in {1,3,4,7,8,9,11,12,14,15,16} {
        \fill[cgrey!45] ({\k-1},{16-\k}) rectangle ({\k},{17-\k});
    }
    \draw[thick] (0,0) rectangle (16,16);
    \node[below] at (8,0) {\small $\Heff^{\mathrm{emb}}$, $\nmax = 16$};
    \begin{scope}[xshift=7.5cm]
        \fill[cblue!60] (0,10) rectangle (0.6,10.6); \node[right] at (0.8,10.3) {\small entries of $\Heff$ ($\Nx = \{2,5,6,10,13\}$)};
        \fill[cgrey!45] (0,8.6) rectangle (0.6,9.2); \node[right] at (0.8,8.9) {\small identity on inactive indices};
        \node[align=left, right] at (0,5.5) {\small A column entering $\Nx$ fills its row and\\ \small column (row addition); a column leaving\\ \small $\Nx$ empties them (row deletion).};
    \end{scope}
\end{tikzpicture}
\caption{Fixed-length representation of the effective reduced Hessian: the active indices
$\Nx$ are scattered over the fixed index range $1,\dots,\nmax$, and every other index carries
an identity row.}
\label{fig:fixed-length}
\end{figure}

The identity rows do not interfere with the computation of the search direction:

\begin{lemma}[Decoupling of inactive indices] \label{lem:embedding}
Let $\Heff^{\mathrm{emb}} = LDL^\top$ be the $LDL^\top$ factorization without pivoting. Then
for every $k\notin\Nx$, row and column $k$ of $L$ are the unit vectors $e_k^\top$ and $e_k$,
and $D_{kk} = 1$. Consequently, for a right-hand side $r$ that vanishes outside $\Nx$, the
solution of $\Heff^{\mathrm{emb}}\eta = r$ vanishes outside $\Nx$, and its restriction to $\Nx$
solves $\Heff\eta_x = r_{\Nx}$.
\end{lemma}

\begin{proof}
Let $P$ be the permutation that moves the indices of $\Nx$ to the front. Then
$P^\top\Heff^{\mathrm{emb}}P = \diag(\Heff, I)$, because an inactive index has a $1$ on the
diagonal and zeros elsewhere in its row and column. With $\Heff = L_1D_1L_1^\top$, this
block-diagonal matrix has the block-diagonal factorization
\[
\diag(\Heff, I) = \diag(L_1,I)\,\diag(D_1,I)\,\diag(L_1,I)^\top,
\]
and the $LDL^\top$ factorization without pivoting is unique. Applying the
permutation backwards only moves rows and columns of these factors, so in the natural order
every inactive index $k$ still has row and column $e_k$ in $L$ and $D_{kk} = 1$. For the second
statement, $\diag(\Heff, I)\,[\eta_x;\eta_0] = [r_{\Nx}; 0]$ gives $\Heff\eta_x = r_{\Nx}$ and
$\eta_0 = 0$.
\end{proof}


\section{Expanding the Working Set}
\label{sec:expanding}

A step is blocked when a free variable reaches zero. The blocking index $q$ is then added to
the working set. There are three cases, depending on the set that $q$ belongs to. For a block
$k$ we write
\[
\mathcal{N}_x^{(k)} = \{\, j\in\Nx : \psi(j) = k \,\}
\]
for the nullspace columns of block $k$; these are the columns whose nullspace vectors
$z_j = e_j - e_{S_k}$ refer to the representative $S_k$.

\begin{lemma}[Ordering after an expansion] \label{lem:expand-ordering}
Let $q$ be the blocking index. The extended ordering is updated as follows.
\begin{enumerate}
    \item $q\in\Nx$: $\;\Nx\gets\Nx\setminus\{q\}$, $\;\Wx\gets\Wx\cup\{q\}$.
    \item $q = S_k\in\Sset$: choose a \textbf{replacement} $r\in\mathcal{N}_x^{(k)}$;
          then $S_k\gets r$, $\;\Nx\gets\Nx\setminus\{r\}$, $\;\Wx\gets\Wx\cup\{q\}$.
    \item $q$ is the free surplus $i\in\Ns$: $\;\Ns\gets\Ns\setminus\{i\}$, $\;\Ws\gets\Ws\cup\{i\}$.
\end{enumerate}
\end{lemma}

In case~(ii) the representative itself reaches zero, and another free column of its block
has to take its place in the list $\Sset$. The next lemma shows that one is always
available.

\begin{lemma}[Existence of a replacement] \label{lem:replacement}
In case~(ii), the set $\mathcal{N}_x^{(k)}$ is not empty.
\end{lemma}

\begin{proof}
By~\eqref{eq:direction-components}, the direction changes the representative by
$p_{S_k} = -\sum_{j\in\mathcal{N}_x^{(k)}}\eta_j$. The representative can only reach zero if it
decreases, $p_{S_k}<0$, so the sum is not empty.
\end{proof}

Any $r\in\mathcal{N}_x^{(k)}$ is a valid replacement. A good choice is the column with the largest
weight, $r = \arg\max_{j\in\mathcal{N}_x^{(k)}}x_j$, since it is the least likely to block in the
next step.

\begin{theorem}[Update of the nullspace basis] \label{thm:Z-update}
After an expansion, the nullspace basis changes as follows.
\begin{enumerate}
    \item $q\in\Nx$: the vector $z_q$ is removed.
    \item $q = S_k$ with replacement $r$: the vector $z_r$ is removed, and for every
          $j\in\mathcal{N}_x^{(k)}\setminus\{r\}$ the vector $z_j$ is replaced by
          \[
          z_j' = e_j - e_r = z_j - z_r .
          \]
          All other vectors are unchanged.
    \item $q$ the surplus $i$: the vector $e_{n+i}$ is removed.
\end{enumerate}
\end{theorem}

\begin{proof}
By \cref{cor:z-explicit} there is one basis vector per index of $\Nx$ and of $\Ns$, given by
$z_j = e_j - e_{\rep(j)}$ and $z_i = e_{n+i}$. In cases~(i) and~(iii) one index leaves $\Nx$ or
$\Ns$, and the representatives do not change, so only its vector disappears. In case~(ii), $r$
leaves $\Nx$ and the representative of block $k$ changes from $q$ to $r$. For
$j\in\mathcal{N}_x^{(k)}\setminus\{r\}$ the new vector is
$z_j' = e_j - e_r = (e_j - e_q) - (e_r - e_q) = z_j - z_r$. The vectors of other blocks refer
to other representatives and do not change.
\end{proof}

\begin{theorem}[Update of the effective reduced Hessian, expansion] \label{thm:H-expand}
After an expansion, $\Heff$ changes as follows.
\begin{enumerate}
    \item $q\in\Nx$: row and column $q$ are removed.
    \item $q = S_k$ with replacement $r$: row and column $r$ are removed, and for
          $i,j\in\mathcal{N}_x^{(k)}\setminus\{r\}$ and $\ell\in\Nx\setminus\mathcal{N}_x^{(k)}$
          \[
          (\Heff')_{\ell j} = (\Heff)_{\ell j} - (\Heff)_{\ell r},\qquad
          (\Heff')_{ij} = (\Heff)_{ij} - (\Heff)_{ir} - (\Heff)_{rj} + (\Heff)_{rr}.
          \]
    \item $q$ the surplus of inequality row $i$: row $i$ of $A_I$ joins $\Aeff$, and
          \begin{equation} \label{eq:rank-one-up}
          \Heff' = \Heff + v_iv_i^\top,\qquad
          v_i = Z_x^\top (A_I)_{i,:}^\top,\qquad (v_i)_j = (A_I)_{ij} - (A_I)_{i,\rep(j)}\;\;(j\in\Nx).
          \end{equation}
\end{enumerate}
\end{theorem}

\begin{proof}
We use $(\Heff)_{\ell j} = (\Aeff z_\ell)^\top(\Aeff z_j)$ from \cref{cor:Aeff}.

(i) The vector $z_q$ disappears, so the column $\Aeff z_q$ of $\Aeff Z_x$ is removed, and with
it row and column $q$ of the Gram matrix $\Heff$.

(ii) The vector $z_r$ disappears, which removes row and column $r$. For the remaining columns
of block $k$, $\Aeff z_j' = \Aeff z_j - \Aeff z_r$ by \cref{thm:Z-update}. For
$\ell\notin\mathcal{N}_x^{(k)}$ the vector $z_\ell$ is unchanged, so
$(\Heff')_{\ell j} = (\Aeff z_\ell)^\top(\Aeff z_j - \Aeff z_r) = (\Heff)_{\ell j} - (\Heff)_{\ell r}$.
For two columns of block $k$, expanding the product
$(\Aeff z_i - \Aeff z_r)^\top(\Aeff z_j - \Aeff z_r)$ gives the second formula.

(iii) The basis $Z_x$ does not change, but $\Aeff$ gains the row $(A_I)_{i,:}$. Hence
\[
\Aeff' Z_x = \begin{bmatrix}\Aeff Z_x\\ (A_I)_{i,:}Z_x\end{bmatrix},\qquad
\Heff' = (\Aeff Z_x)^\top(\Aeff Z_x) + \bigl((A_I)_{i,:}Z_x\bigr)^\top\bigl((A_I)_{i,:}Z_x\bigr)
= \Heff + v_iv_i^\top .
\]
Entry $j$ of $v_i = Z_x^\top(A_I)_{i,:}^\top$ is the product of column $j$ of $Z_x$, which is
$e_j - e_{\rep(j)}$, with row $i$ of $A_I$, i.e.\ $(v_i)_j = (A_I)_{ij} - (A_I)_{i,\rep(j)}$.
\end{proof}

The vector $v_i$ is row $i$ of $B_I = A_IZ_x$. Its entry $(v_i)_j$ is nonzero only if exactly
one of the columns $j$ and $\rep(j)$ has a nonzero coefficient in row $i$; in crew scheduling,
if exactly one of the two schedules covers flight $i$.

\begin{prop} \label{prop:expand-pd}
If $\Heff$ is positive definite before an expansion, it is positive definite after it.
\end{prop}

\begin{proof}
By \cref{cor:dimension}, $\Heff$ is positive definite if and only if the vectors $\Aeff z_j$,
$j\in\Nx$, are linearly independent.
(i) Removing one vector from a linearly independent set leaves it independent.
(ii) The new set consists of the unchanged vectors $\Aeff z_\ell$ and the vectors
$\Aeff z_j - \Aeff z_r$ for $j\in\mathcal{N}_x^{(k)}\setminus\{r\}$. A vanishing linear combination
$\sum_\ell\alpha_\ell\Aeff z_\ell + \sum_j\beta_j(\Aeff z_j - \Aeff z_r) = 0$ is a vanishing
combination of the old, independent vectors, in which the coefficient of $\Aeff z_j$ is
$\beta_j$ and that of $\Aeff z_\ell$ is $\alpha_\ell$; hence all $\alpha_\ell$ and $\beta_j$ are zero.
(iii) For every $\eta_x\ne0$, $\eta_x^\top(\Heff + v_iv_i^\top)\eta_x = \eta_x^\top\Heff\eta_x +
(v_i^\top\eta_x)^2 > 0$.
\end{proof}


\section{Reducing the Working Set}
\label{sec:reducing}

At a subspace minimizer with a negative multiplier, the index with the most negative
multiplier is released from the working set. By \cref{cor:multipliers-minimizer}, the
multipliers of columns (reduced costs) and of surpluses (duals) are both measured in units
of cost, so they can be compared directly. Releasing an index adds a basis vector to $Z$:

\begin{theorem}[Update of the effective reduced Hessian, reduction] \label{thm:H-reduce}
The two kinds of reduction change $\Heff$ as follows.
\begin{enumerate}
    \item[(iv)] Column $r\in\Wx$ is released ($\Wx\gets\Wx\setminus\{r\}$, $\Nx\gets\Nx\cup\{r\}$).
          The basis gains $z_r = e_r - e_{\rep(r)}$, and with $u = \Aeff z_r$,
          \begin{equation} \label{eq:bordered}
          \Heff' = \begin{bmatrix}\Heff & h\\ h^\top & h_{rr}\end{bmatrix},\qquad
          h = (\Aeff Z_x)^\top u,\qquad h_{rr} = u^\top u .
          \end{equation}
    \item[(v)] The surplus of inequality row $i$ is released ($\Ws\gets\Ws\setminus\{i\}$,
          $\Ns\gets\Ns\cup\{i\}$). The basis gains $e_{n+i}$, row $i$ leaves $\Aeff$, and
          \begin{equation} \label{eq:rank-one-down}
          \Heff' = \Heff - v_iv_i^\top
          \end{equation}
          with $v_i = Z_x^\top(A_I)_{i,:}^\top$ as in~\eqref{eq:rank-one-up}.
\end{enumerate}
\end{theorem}

\begin{proof}
(iv) The new column part of the basis is $Z_x' = [\,Z_x\;\;z_r\,]$, so
$\Aeff Z_x' = [\,\Aeff Z_x\;\;u\,]$ and
\[
\Heff' = \begin{bmatrix}(\Aeff Z_x)^\top\\ u^\top\end{bmatrix}\begin{bmatrix}\Aeff Z_x & u\end{bmatrix}
= \begin{bmatrix}\Heff & (\Aeff Z_x)^\top u\\ u^\top\Aeff Z_x & u^\top u\end{bmatrix}.
\]
(v) $Z_x$ does not change, and $\Aeff$ loses the row $(A_I)_{i,:}$. This is the reverse of
case~(iii) of \cref{thm:H-expand}, so $\Heff = \Heff' + v_iv_i^\top$.
\end{proof}

\begin{lemma}[Descent after a reduction] \label{lem:descent}
Suppose the reduced Hessian is positive definite after the reduction. Then the next search
direction moves the released variable away from its bound: $p_r>0$ in case~(iv), and
$(p_s)_i>0$ in case~(v).
\end{lemma}

\begin{proof}
Before the reduction the iterate is a subspace minimizer, so every entry of the reduced
gradient $Z^\top g$ is zero. The reduction adds one basis vector $z$ ($z_r$ or $e_{n+i}$) and
leaves the others unchanged, so the new reduced gradient has the single nonzero entry
$z^\top g$. For $z = z_r$ it is $g_r - g_{\rep(r)}$, and for $z = e_{n+i}$ it is $\pi_i$; by
\cref{cor:multipliers-minimizer} both equal the multiplier $\nu<0$ of the released index. Let
$e$ denote the unit vector of the new entry. The new reduced system
$\Hbar'\eta = -\delta Z'^\top g = -\delta\nu\,e$ gives $\eta = -\delta\nu\,\Hbar'^{-1}e$, whose own
component is
\[
\eta_e = -\delta\nu\,(\Hbar'^{-1})_{ee} > 0,
\]
since $\Hbar'^{-1}$ is positive definite and $\nu<0$. By~\eqref{eq:direction-components} this
component is the change of the released variable ($p_r = \eta_r$, $(p_s)_i = \eta_i$).
\end{proof}

Unlike an expansion, a reduction can make the reduced Hessian singular: the new column $u$ in
case~(iv) can lie in the span of the old ones, and the downdate in case~(v) can remove the
last direction of curvature. This is the subject of \cref{chap:singular}.


\section{Updating the $LDL^\top$ Factorization}
\label{sec:ldl-updates}

The updates of \cref{sec:expanding,sec:reducing} are applied to the factorization
$\Heff^{\mathrm{emb}} = LDL^\top$, with $L$ unit lower triangular and $D$ diagonal, in the
natural order of the indices. We use the sparse direct solver
CHOLMOD~\cite{chen2008cholmod}, which provides the three required modifications
\cite{davis1999modifying, davis2005row}:

\begin{center}
\renewcommand{\arraystretch}{1.25}
\begin{tabular}{lll}
\toprule
change of the working set & change of $\Heff^{\mathrm{emb}}$ & CHOLMOD routine \\
\midrule
(i) column leaves $\Nx$ & row/column $q$ becomes $e_q$ & \code{rowdel} \\
(ii) representative swap & rows/columns of $\mathcal{N}_x^{(k)}$ rewritten & \code{rowdel} + \code{rowadd} \\
(iii) surplus becomes active & $+\,v_iv_i^\top$ & \code{updown} (update) \\
(iv) column enters $\Nx$ & row/column $r$ set to $[h; h_{rr}]$ & \code{rowadd} \\
(v) surplus becomes free & $-\,v_iv_i^\top$ & \code{updown} (downdate) \\
\bottomrule
\end{tabular}
\end{center}

\subsection{Row additions and deletions}

The following proposition shows what a row addition computes. It explains the cost of the
operation and the meaning of the new pivot, which is used to detect singularity in
\cref{chap:singular}.

\begin{prop}[Row addition] \label{prop:rowadd}
Let row and column $k$ of $\Heff^{\mathrm{emb}}$ currently be $e_k$, and partition the indices
into $1 = \{1,\dots,k-1\}$, $k$ and $3 = \{k+1,\dots,\nmax\}$:
\[
L = \begin{bmatrix} L_{11} & & \\ 0 & 1 & \\ L_{31} & 0 & L_{33}\end{bmatrix},\qquad
D = \diag(D_1, 1, D_3).
\]
If row and column $k$ are set to $(h_1, h_{kk}, h_3)$, the new factor has
\begin{align*}
\ell_{k1} &= D_1^{-1}L_{11}^{-1}h_1 &&\text{(triangular solve with the leading block)},\\
d_k &= h_{kk} - \ell_{k1}^\top D_1\ell_{k1} &&\text{(new pivot)},\\
\ell_{3k} &= \bigl(h_3 - L_{31}D_1\ell_{k1}\bigr)/d_k &&\text{(new column below $k$)},\\
L_{33}'D_3'L_{33}'^\top &= L_{33}D_3L_{33}^\top - d_k\,\ell_{3k}\ell_{3k}^\top &&\text{(rank-one downdate)},
\end{align*}
while $L_{11}$, $D_1$ and $L_{31}$ are unchanged.
\end{prop}

\begin{proof}
The new factor has the form
$L' = \left[\begin{smallmatrix} L_{11} & & \\ \ell_{k1}^\top & 1 & \\ L_{31} & \ell_{3k} & L_{33}'\end{smallmatrix}\right]$,
$D' = \diag(D_1, d_k, D_3')$. Multiplying out $L'D'L'^\top$ block by block and comparing with
the new matrix:
\begin{itemize}
    \item block $(k,1)$: $\ell_{k1}^\top D_1L_{11}^\top = h_1^\top$, which gives $\ell_{k1}$;
    \item entry $(k,k)$: $\ell_{k1}^\top D_1\ell_{k1} + d_k = h_{kk}$, which gives $d_k$;
    \item block $(3,k)$: $L_{31}D_1\ell_{k1} + \ell_{3k}d_k = h_3$, which gives $\ell_{3k}$;
    \item block $(3,3)$: $L_{31}D_1L_{31}^\top + d_k\ell_{3k}\ell_{3k}^\top + L_{33}'D_3'L_{33}'^\top$ must
          equal the unchanged block $L_{31}D_1L_{31}^\top + L_{33}D_3L_{33}^\top$, which gives the
          downdate.
\end{itemize}
The blocks $(1,1)$ and $(3,1)$ do not involve row $k$ and stay as they are.
\end{proof}

A row deletion is the reverse operation: row and column $k$ are reset to $e_k$, which
requires the rank-one \emph{update} $L_{33}D_3L_{33}^\top + d_k\ell_{3k}\ell_{3k}^\top$ of the trailing
block. Two consequences of \cref{prop:rowadd} are important.

\begin{itemize}
    \item If $k$ is larger than every active index, the trailing block contains no active
          index, the downdate is trivial, and a row addition costs a single triangular solve
          with the active part, $O(|\Nx|^2)$ for a dense factor. This is used by the
          refactorization in \cref{sec:refactorization}.
    \item The determinant of the new active part is $\det\Heff' = \det\Heff\cdot d_k$, since
          the determinant of an $LDL^\top$ factorization is the product of the pivots and the
          other pivots of the active part do not change. A vanishing pivot $d_k$ therefore
          signals a singular $\Heff'$ (\cref{chap:singular}). If $k$ is the last active index,
          $d_k = h_{kk} - h^\top\Heff^{-1}h$ in the notation of~\eqref{eq:bordered}.
\end{itemize}

\subsection{Rank-one modifications}

The surplus cases (iii) and (v) modify the whole matrix by $\pm v_iv_i^\top$. The factorization
$LDL^\top\pm vv^\top = L'D'L'^\top$ is computed by the classical rank-one update and downdate
algorithms \cite{gill1974methods, davis1999modifying}, which process the columns of $L$ along
the path in the elimination tree that starts at the first nonzero of $v$. Their cost is
proportional to the number of entries of $L$ in these columns. An update keeps the matrix
positive definite (\cref{prop:expand-pd}). A downdate can make it singular; by the matrix
determinant lemma,
\[
\det\bigl(\Heff - v_iv_i^\top\bigr) = \det\Heff\cdot\bigl(1 - v_i^\top\Heff^{-1}v_i\bigr),
\]
so this happens exactly when $\omega_i := 1 - v_i^\top\Heff^{-1}v_i$ vanishes.
\Cref{chap:singular} shows that $\omega_i\ge0$ always holds for the downdates of the
active-set method, and treats the case $\omega_i = 0$.

\subsection{Sparse solves}

The right-hand side of~\eqref{eq:eff-system} vanishes outside $\Nx$. By \cref{lem:embedding},
the solve only has to touch the rows of $L$ that can be reached from $\Nx$ in the graph of
$L$, the set $\reach_L(\Nx)$. CHOLMOD supports such solves directly. Their cost is independent
of $\nmax$, which is what makes a large capacity $\nmax$ affordable.

\subsection{Assembly of a new row}

A row addition in case~(iv) needs the vector $h$ of~\eqref{eq:bordered}. By~\eqref{eq:implicit},
$h_k = u_k^\top M_\rho u$ with $u = a_r - a_{\rep(r)}$. In an implementation, $M_\rho u$ is
scattered into a dense $m$-vector once, with the entries of the rows of free surpluses set to
zero. Then $h_k = a_k^\top(M_\rho u) - a_{\rep(k)}^\top(M_\rho u)$ is computed for every
$k\in\Nx$, where the products with the $K$ representatives are computed once and reused. The
cost is $O\bigl((|\Nx| + K)\,\bar\ell\bigr)$, independent of $n$. A representative swap
(case~(ii)) rewrites the rows of $\mathcal{N}_x^{(k)}$ by a row deletion followed by a row addition
of the freshly assembled row; this costs the same as applying the formulas of
\cref{thm:H-expand}, and it needs no stored copy of $\Heff$.


\section{Structural Zeros and Periodic Refactorization}
\label{sec:refactorization}

A row deletion sets row and column $k$ of $L$ to $e_k$ \emph{numerically}, but a sparse
factorization keeps the positions of the removed entries in its sparsity pattern: the
pattern of a column of a simplicial factor can only grow. The same holds for rank-one
modifications. Over a long sequence of updates, the stored pattern of $L$ therefore becomes
the union of the patterns of all factors that ever existed. To quantify this effect we
distinguish the following quantities, which are measured in \cref{sec:results-fill}.

\begin{definition}[Stored and effective entries] \label{def:stored-effective}
The \textbf{stored entries} of $L$ are the off-diagonal positions in its sparsity pattern.
The \textbf{effective entries} are the stored entries whose value is nonzero. The difference
are the \textbf{structural zeros}. The \textbf{dense equivalent} is $|\Nx|(|\Nx|-1)/2$, the
number of off-diagonal entries of a dense factor of $\Heff$.
\end{definition}

Every solve and every update iterates over the stored entries of the columns it touches,
including the structural zeros. Moreover, structural zeros can connect active indices to
indices that are no longer active, which enlarges $\reach_L(\Nx)$. The stored pattern only
returns to the effective one when the factorization is rebuilt.

\begin{definition}[Refactorization] \label{def:refactorization}
A \textbf{refactorization} resets the factor to the identity and re-activates all indices of
$\Nx$ by row additions \emph{in increasing index order}, with the rows assembled as in
\cref{sec:ldl-updates}.
\end{definition}

\begin{prop}[Cost of a refactorization] \label{prop:refactor-cost}
A refactorization costs $O(|\Nx|^3)$ arithmetic operations for the factor (about
$|\Nx|^3/3$ for a dense $\Heff$) and $O\bigl(|\Nx|(|\Nx|+K)\,\bar\ell\bigr)$ for the assembly of
the rows. Its stored pattern contains no structural zeros.
\end{prop}

\begin{proof}
When the indices are added in increasing order, every new index is the largest active index,
so by \cref{prop:rowadd} each row addition is a single triangular solve with the current
active part, which costs $O(t^2)$ for the $t$-th index. Summing over $t = 1,\dots,|\Nx|$ gives
$O(|\Nx|^3)$, about $|\Nx|^3/3$ as for a dense Cholesky factorization. Each row costs
$O\bigl((|\Nx|+K)\,\bar\ell\bigr)$ to assemble (\cref{sec:ldl-updates}). Since no row is ever
deleted during a refactorization, every stored entry is computed from the current matrix.
\end{proof}

The method refactorizes after every $\tau$ updates of the factorization (row additions, row
deletions and rank-one modifications), and whenever a pivot has become too small
(\cref{chap:singular}). The \textbf{refactorization interval} $\tau$ balances the cost of a
rebuild against the growth of the structural zeros, which is roughly proportional to the number
of updates. Its effect is studied in \cref{sec:results-fill}.

\begin{remark}[Natural ordering] \label{rem:natural-ordering}
Sparse direct solvers usually reorder a matrix to reduce fill-in. For $\Heff$ this would gain
nothing: the effective factor is almost exactly as dense as a dense factor of $\Heff$
(\cref{sec:results-fill}), because columns of the same block share their representative and
are therefore strongly coupled. The natural ordering has two advantages: index $j$ of the
factor is column $j$ of the master problem, and the refactorization of
\cref{def:refactorization} only appends rows.
\end{remark}


\section{Hot Start Between Column Generation Iterations}
\label{sec:hot-start}

Between two master problems, regularized column generation adds columns and moves the
stability center (\cref{alg:rcg}). Both changes can be absorbed without touching the
factorization.

\begin{theorem}[Hot start] \label{thm:hot-start}
Let $y = [x;s]$ be the final iterate of a master problem, with its extended ordering and the
factorization of $\Heff$.
\begin{enumerate}
    \item \textbf{New column.} If a column $n+1\le\nmax$ is added with weight $x_{n+1} = 0$
          and placed in the working set $\Wx$, then $y$ remains feasible, $\pi(y)$ is
          unchanged, and the nullspace basis, $\Aeff$, $\Heff$ and its factorization are
          unchanged. The multiplier of the new column is its reduced cost, which is negative
          because pricing selected it.
    \item \textbf{New center.} If the center changes from $\pihat$ to $\pihat + \Delta$,
          then $y$ remains feasible, the ordering and the factorization are unchanged, and the
          dual estimate becomes $\pi(y) + \Delta$.
\end{enumerate}
\end{theorem}

\begin{proof}
(i) The new column has weight zero, so the convexity constraint of its block and the product
$\Abar y$ are unchanged; hence $y$ stays feasible and $\pi(y) = \pihat + \frac{1}{\delta}(b - \Abar y)$
is unchanged. The new index is placed in $\Wx$, so the sets $\Sset$, $\Nx$, $\Ns$, and with them
$Z$, $\Aeff$ and $\Heff$, are unchanged. In the fixed-length representation, index $n+1$
already carries an identity row, so the factorization is unchanged too. By
\cref{cor:multipliers-minimizer}, the multiplier of a working-set column is its reduced cost
$c_{n+1} - a_{n+1}^\top\pi - \mu_{\psi(n+1)}$, which pricing found to be negative.
(ii) The constraints $\Cbar y = \one$, $y\ge0$ do not involve $\pihat$, and
$\Heff = (\Aeff Z_x)^\top(\Aeff Z_x)$ does not depend on it either. Finally
$\pi(y) = (\pihat + \Delta) + \frac{1}{\delta}(b - \Abar y)$ is the old dual estimate plus $\Delta$.
\end{proof}

After a new column, the previous optimum is a subspace minimizer whose only negative
multipliers are those of the new columns. The next master problem therefore typically starts
by releasing the new column (case~(iv)). After a new center, the gradient changes and the
iterate is in general no longer a subspace minimizer; the next master problem starts with a
direction computation on the inherited factorization. The surpluses are never added or
removed, and their status is carried over like that of the columns.

\subsection{Center-only iterations}

The hot start also explains the behaviour of the center-only iterations of
\cref{sec:center-update}, in which the center moves to the current dual and no column is
added. Let $\Pi = \Abar Z\,\Hbar^{-1}(\Abar Z)^\top$ denote the orthogonal projector onto the range
of $\Abar Z$.

\begin{prop}[Center-only iterations] \label{prop:center-only}
Let $y$ be optimal for the current master problem, with residual $r = b - \Abar y$ and dual
$\pi = \pihat + r/\delta$, and let the center move to $\pihat' = \pi$. Suppose that the next
master problem takes one full step and keeps the working set. Then
\begin{enumerate}
    \item the new residual is $r' = (I-\Pi)\,r$ and the new dual is $\pi' = \pi + (I-\Pi)\,r/\delta$;
    \item if the working set stays the same in the following center-only iterations, the
          primal iterate no longer changes, and every further iteration moves the dual by the
          fixed vector $r'/\delta$:
          \[
          \pi^{(t+1)} = \pi^{(t)} + \frac{r'}{\delta},\qquad
          (\nu_x)_j^{(t+1)} = (\nu_x)_j^{(t)} - \frac{u_j^\top r'}{\delta},\qquad
          (\nu_s)_i^{(t+1)} = (\nu_s)_i^{(t)} + \frac{r'_i}{\delta}.
          \]
\end{enumerate}
\end{prop}

\begin{proof}
(i) The center moves by $\Delta = \pihat' - \pihat = \pi - \pihat = r/\delta$. By
\cref{thm:hot-start}, the dual at the old iterate becomes $\pi + \Delta$, and since the gradient
is $g = \cbar - \Abar^\top\pi$ (\cref{lem:gradient-dual}), it changes by $-\Abar^\top\Delta$. The
old iterate was optimal, so its reduced gradient was zero, and the new reduced gradient is
\[
Z^\top g' = Z^\top g - Z^\top\Abar^\top\Delta = -(\Abar Z)^\top\Delta .
\]
The reduced system $\Hbar\eta = -\delta Z^\top g'$ gives $\eta = \delta\,\Hbar^{-1}(\Abar Z)^\top\Delta$,
hence $\Abar p = \Abar Z\eta = \delta\,\Pi\Delta$. After the full step, the dual is
\[
\pi' = (\pi + \Delta) - \tfrac{1}{\delta}\Abar p = \pi + \Delta - \Pi\Delta = \pi + (I - \Pi)\Delta,
\]
and by \cref{prop:residual} the new residual is $r' = \delta(\pi' - \pihat') = \delta(I-\Pi)\Delta = (I-\Pi)\,r$.

(ii) The vector $r' = (I-\Pi)r$ is orthogonal to the range of $\Abar Z$, i.e.\
$(\Abar Z)^\top r' = 0$. In the next center-only iteration the center moves by
$\Delta' = r'/\delta$, and the same computation as in~(i) gives the reduced gradient
$-(\Abar Z)^\top\Delta' = 0$: the direction is zero and $y$ does not move. The dual becomes
$\pi' + \Delta'$; since the new center is $\pi'$, the residual is
$\delta\bigl((\pi'+\Delta') - \pi'\bigr) = r'$ again, and the argument repeats. The multipliers follow from \cref{cor:multipliers-minimizer}. A shift of $\pi$ by
$\Delta'$ changes $g_j$ by $-a_j^\top\Delta'$, hence $(\nu_x)_j = g_j - g_{\rep(j)}$ by
$-u_j^\top\Delta'$; and $(\nu_s)_i = \pi_i$ changes by $\Delta'_i$.
\end{proof}

\Cref{prop:center-only} shows that in a center-only phase the dual walks along a fixed
direction with step length $\lVert r'\rVert/\delta$, while the primal iterate stays put. The
walk ends when a multiplier changes sign, so that the working set changes, or when pricing
finds an improving column. If the residual $r'$ is small, many such iterations are needed,
which explains the long tail of center-only iterations observed in \cref{sec:results-phases}.
Each of them is cheap: one direction computation, which returns zero, and one multiplier test.


\section{The Complete Method}
\label{sec:complete}

\Cref{alg:hot-qpas} combines the results of
\cref{chap:qpas,chap:nullspace,chap:updates}. It is called in the first line of the loop of
regularized column generation (\cref{alg:rcg}), and its state, consisting of the iterate, the
extended ordering and the factorization, is kept from one call to the next.

\begin{algobox}{Hot-Started Nullspace Active-Set Method}{alg:hot-qpas}
\State \textbf{Input:} state of the previous call (iterate $y$, ordering, factor of $\Heff$),
       new columns, new center $\pihat$
\State Put the new columns into $\Wx$ with weight $0$ \Comment{\cref{thm:hot-start}}
\State $\pi\gets\pihat + (b - \Abar y)/\delta$
\Repeat
    \If{the last step was not a full step}
        \State Compute $p$ by \cref{alg:nullspace-iteration} (lines 2--4)
    \Else
        \State $p\gets0$ \Comment{\cref{prop:full-step}}
    \EndIf
    \If{$p\approx0$}
        \State Compute the multipliers (\cref{cor:multipliers-minimizer})
        \If{all multipliers $\ge-\varepsilon$} \textbf{return} $y$, $\pi$, $\mu$ \EndIf
        \State Release the most negative index: \code{rowadd} (column) or downdate (surplus)
        \State If the new pivot is too small, resolve the singularity (\cref{alg:singular})
    \Else
        \State Ratio test over $\Sset\cup\Nx\cup\Ns$; $y\gets y + \alpha p$;
               $\pi\gets\pi - \frac{\alpha}{\delta}\Abar p$
        \If{blocked by $q$}
            \State Expand: \code{rowdel} ($q\in\Nx$), representative swap ($q\in\Sset$), or
                   update ($q\in\Ns$)
        \EndIf
    \EndIf
    \State Refactorize if $\tau$ updates have passed since the last refactorization
\Until{returned}
\end{algobox}
